\section{Snap Targets and Normalization}
\label{app:snap}

After training, learned parameters are continuous floating-point values.
To recover closed-form symbolic expressions, each parameter is rounded
(``snapped'') to the nearest value in a predefined table of
mathematically meaningful targets. Table~\ref{tab:snap-targets}
enumerates the complete set of snap targets used throughout this work.

\begin{table}[h]
\centering
\caption{Complete table of snap targets used for rounding learned parameters to clean symbolic values.}
\label{tab:snap-targets}
\begin{tabular}{lrl}
\toprule
\textbf{Target Value} & \textbf{Numeric Value} & \textbf{Label} \\
\midrule
$0$             &  0.000000 & zero \\
$+1$            &  1.000000 & one \\
$-1$            & $-1.000000$ & neg one \\
$+2$            &  2.000000 & two \\
$-2$            & $-2.000000$ & neg two \\
$+3$            &  3.000000 & three \\
$-3$            & $-3.000000$ & neg three \\
$+\tfrac{1}{2}$ &  0.500000 & half \\
$-\tfrac{1}{2}$ & $-0.500000$ & neg half \\
$+\tfrac{1}{3}$ &  0.333333 & third \\
$-\tfrac{1}{3}$ & $-0.333333$ & neg third \\
$+\tfrac{2}{3}$ &  0.666667 & two-thirds \\
$-\tfrac{2}{3}$ & $-0.666667$ & neg two-thirds \\
$+\tfrac{1}{4}$ &  0.250000 & quarter \\
$-\tfrac{1}{4}$ & $-0.250000$ & neg quarter \\
$+\tfrac{1}{5}$ &  0.200000 & fifth \\
$-\tfrac{1}{5}$ & $-0.200000$ & neg fifth \\
$+\tfrac{3}{2}$ &  1.500000 & three-halves \\
$-\tfrac{3}{2}$ & $-1.500000$ & neg three-halves \\
$+\pi$          &  3.141593 & pi \\
$-\pi$          & $-3.141593$ & neg pi \\
$+\tfrac{\pi}{2}$ &  1.570796 & pi/2 \\
$-\tfrac{\pi}{2}$ & $-1.570796$ & neg pi/2 \\
$+\tfrac{\pi}{4}$ &  0.785398 & pi/4 \\
$-\tfrac{\pi}{4}$ & $-0.785398$ & neg pi/4 \\
$+e$            &  2.718282 & euler \\
$-e$            & $-2.718282$ & neg euler \\
$+\sqrt{2}$    &  1.414214 & sqrt2 \\
$-\sqrt{2}$    & $-1.414214$ & neg sqrt2 \\
$+\ln 2$       &  0.693147 & ln2 \\
$-\ln 2$       & $-0.693147$ & neg ln2 \\
\bottomrule
\end{tabular}
\end{table}

Given a learned parameter $\hat{\theta}$, snapping selects
$\theta^* = \arg\min_{\theta \in \mathcal{T}} |\hat{\theta} - \theta|$,
where $\mathcal{T}$ is the set of targets above, subject to the
constraint that $|\hat{\theta} - \theta^*| < \epsilon_{\mathrm{snap}}$.
If no target lies within the tolerance $\epsilon_{\mathrm{snap}}$, the
raw floating-point value is retained.

\subsection{Normalization of \texttt{SeparableTerm}}

The parameterization of a \texttt{SeparableTerm} as a product
\begin{equation}
  c \;\prod_k f_k(a_k x_k + b_k),
  \label{eq:separable-form}
\end{equation}
where each $f_k \in \{\exp, \sin, \cos, \mathrm{id}\}$, possesses
\emph{gauge freedom}: multiple distinct parameter settings yield the
same function. For example,
\begin{equation}
  c \cdot \exp(a x + b) \;=\; \bigl[c \cdot e^{b}\bigr] \cdot \exp(a x),
\end{equation}
so the coefficient $c$ and bias $b$ can trade off arbitrarily without
changing the represented function. If left unnormalized, the optimizer
may converge to any point along this gauge orbit, placing the learned
parameters far from the snap targets and causing snapping to produce
incorrect symbolic expressions. Normalization breaks this symmetry by
imposing canonical constraints before snapping is applied.

The normalization algorithm proceeds in three steps:

\paragraph{Step 1: Exponential bias absorption.}
For each exponential factor $\exp(a_k x_k + b_k)$ in the term, the
bias is absorbed into the coefficient:
\begin{equation}
  c \leftarrow c \cdot e^{b_k}, \qquad b_k \leftarrow 0.
\end{equation}
After this step, every exponential factor has the form $\exp(a_k x_k)$
and the coefficient $c$ carries the full scale.

\paragraph{Step 2: Trigonometric bias reduction.}
For each sinusoidal factor $\sin(a_k x_k + b_k)$ or
$\cos(a_k x_k + b_k)$, the bias $b_k$ is reduced to the principal
interval $[-\pi, \pi]$ via modular arithmetic:
\begin{equation}
  b_k \leftarrow \bigl((b_k + \pi) \bmod 2\pi\bigr) - \pi.
\end{equation}
If, after reduction, $|b_k| > \tfrac{\pi}{2} + \varepsilon$, a
half-period shift is applied. Using the identity
$\sin(x + \pi) = -\sin(x)$ (and analogously for cosine), we set
\begin{equation}
  b_k \leftarrow b_k - \operatorname{sign}(b_k)\,\pi, \qquad c \leftarrow -c.
\end{equation}
This ensures all trigonometric biases lie in $[-\tfrac{\pi}{2} - \varepsilon,\; \tfrac{\pi}{2} + \varepsilon]$, concentrating them near the
small set of snap targets $\{0, \pm\tfrac{\pi}{4}, \pm\tfrac{\pi}{2}\}$.

\paragraph{Step 3: Sign canonicalization via sine parity.}
If the overall coefficient $c$ is negative and the term contains at
least one sine factor $\sin(a_j x_j + b_j)$, the sign is absorbed
using the odd-symmetry identity $\sin(-z) = -\sin(z)$:
\begin{equation}
  a_j \leftarrow -a_j, \qquad b_j \leftarrow -b_j, \qquad c \leftarrow -c.
\end{equation}
This removes a residual sign ambiguity, ensuring that the coefficient
$c$ is non-negative whenever a sine factor is present.

Together, these three steps place the parameters of each
\texttt{SeparableTerm} into a canonical form in which every degree of
freedom corresponds to a genuinely independent quantity, making the
subsequent snapping procedure both reliable and unambiguous.
